% Note: carrying memory across a code-list revision breaks a precision promise.
% Every number comes from the README files of this repository and of the iyulab-gil package.
% Compile: `tectonic note.tex`, or `pdflatex note.tex` twice (TeX Live 2023 or later).
\documentclass[11pt]{article}
\usepackage[T1]{fontenc}
\usepackage[utf8]{inputenc}
\usepackage[margin=1in]{geometry}
\usepackage{booktabs}
\usepackage{amsmath}
\usepackage[hidelinks]{hyperref}
\usepackage{url}

\title{Carrying memory across a code-list revision breaks a precision promise:\\
two public streams, a watch that notices, and a request for counterexamples}
\author{iyulab}
\date{}

\begin{document}
\maketitle

\begin{abstract}
A system that answers only the records it can answer at a stated precision, and leaves the rest to people, makes a
promise: of the answers it gives, at least this share will be right. We report one way that promise fails and how
quickly the failure can be seen. On two public record streams whose code list was revised so that code numbers were
rewritten (Chicago food inspection violations, 2018 food code; OSHA severe injury reports, OIICS~3.0), a keeper that
carried its memory over the revision promised 0.94 and 0.87 precision for the first 1{,}000 records after it and
realized 0.44 and 0.007. A keeper that started its memory at the revision kept its promise in every range of records in
which it answered. A watch that compares the last 100 answers whose rightness came back against the target raised its
first alarm 54 and 163 records after the revision, with no alarm before it, and none on thirteen other public streams
without a revision. Two score-threshold tools that freeze their threshold before the revision answered 92--100\% of the
records afterwards with fewer than 10\% right. Everything here is two streams, derived (not human-confirmed)
settlements, and one setting; a public kit rebuilds every table from the public data with a published package. We ask
for counterexamples: a stream where carrying memory over such a revision keeps the promise, or where the watch stays
silent after one.
\end{abstract}

\section{Setting}

Many record-keeping systems suggest a field's value from records already settled and let a person confirm or correct
it. Answering only above a threshold is selective prediction: the share answered is the coverage, and one minus the
precision of the answers is the selective risk~\cite{geifman2017}. A threshold chosen on a validation sample for a
stated precision is the usual way to make that risk a promise~\cite{ltt,crc}; its guarantee assumes the records to come
are exchangeable with the sample.

The keeper studied here (Gil, open source~\cite{gil}) does not hold out a sample. It replays the settled records in the
order they were settled, choosing thresholds per layer of evidence (an exact key, similar earlier records) so that the
answers a replay would have given meet a target precision, and it chooses again as settlements arrive. Each answer
carries the promise of the threshold it was given at, and the settled record keeps it, so the realized precision can be
counted against the promise afterwards. Replay in settlement order removes the need for a held-out sample. It does not
remove the exchangeability assumption: a promise chosen on the past is a statement about records like the past.

A code-list revision is the plainest break of that assumption. When a classification edition renumbers its codes, the
same number means something else from one day on, and the records before the day are no longer evidence for the records
after it.

\section{Claims}

\begin{enumerate}
  \item \textbf{Break.} A keeper that carries its memory over a revision that rewrites code numbers promises far more
  than it realizes in the first records after the revision.
  \item \textbf{Split.} A keeper that starts its memory at the revision keeps its promise there.
  \item \textbf{Detection.} A watch on the realized precision of the answers whose rightness came back notices the
  revision within a few hundred records, without being told about it, and does not fire on streams without a revision.
\end{enumerate}

\section{Method}

\paragraph{Streams.} Chicago food inspections (City of Chicago Data Portal): every violation in an inspection is one
record and the decided attribute is the violation number; the revised food code took effect on 2018-07-01 and rewrote
the numbers (1--70 before, 1--64 from it; 45 numbers occur in both). OSHA severe injury reports: every report is one
record and the decided attribute is the nature-of-injury code; from 2024-01 the reports carry OIICS~3.0 codes. The four
most common pre-2024 codes make up 0.73--0.77 of the reports in the last quarter of 2023 and 0.05 in the first quarter
of 2024. A record counts as settled with the code the source lists, in event order --- nobody confirmed or corrected a
suggestion.

\paragraph{Arms.} Both arms are one keeper at a target precision of 0.9 with at least 30 answered records before a
threshold is used. The records after the revision are the same for both. Arm~C is given every record before the
revision with its code unchanged; arm~N is given nothing. Each chooses at the revision, after 250, 500, 1{,}000, 2{,}000
and 4{,}000 records, and from then on whenever the settled records have grown by 10\%. Per range of records after the
revision we report the answered count, the precision of the answers, the promise (mean of what the answering layer
promised) and their difference with a Wilson 95\% interval on the answered records, the promise treated as fixed.

\paragraph{Watch.} A separate run follows each stream without telling the keeper about the revision. The first 5{,}000
records are settled as typed; afterwards each record is answered first, and an answer shown is settled as accepted when
it equals the source code and as corrected when it does not. The watch says the choice no longer stands when, of the
last 100 answers whose rightness came back, fewer than the target less 0.15 (0.75) were right. An alarm is counted at
its first appearance, with at least 5{,}000 records and 100 judged answers since the previous one.

\paragraph{Frozen thresholds.} For comparison, two score-threshold tools --- CalibRoute (\texttt{calibroute-ai}
0.4.0)~\cite{calibroute} and the risk-control module of MAPIE (1.6.0)~\cite{mapie} --- are run as documented with one
off-the-shelf text classifier (multinomial logistic regression on word 1--2-grams of the free text and one-hot
categorical columns; \texttt{scikit-learn} 1.9.1), trained on half of the pre-revision records, calibrated on a quarter
and checked on the last quarter (five seeds), at the same 0.90 target. One variant routes batches of 250 records through
CalibRoute's own shift check; another retrains and refits after the revision at the keeper's choice points.

\section{Results}

\paragraph{Break and split.} Table~\ref{tab:split} gives the first 1{,}000 records after the revision.

\begin{table}[h]
\centering
\caption{First 1{,}000 records after the revision.}
\label{tab:split}
\begin{tabular}{llrrrl}
\toprule
Stream & Arm & Answered & Precision & Promise & Promise $-$ precision (95\%) \\
\midrule
Chicago & N & 356 & 0.975 & 0.969 & $-0.006$ [$-0.018$, $+0.016$] \\
Chicago & C & 655 & 0.441 & 0.940 & $+0.499$ [$+0.461$, $+0.537$] \\
OSHA & N & 0 & -- & -- & -- \\
OSHA & C & 153 & 0.007 & 0.869 & $+0.862$ [$+0.833$, $+0.868$] \\
\bottomrule
\end{tabular}
\end{table}

In the next 4{,}000 records arm~C still promised more than it realized (Chicago 0.802 against 0.611; OSHA 0.536 against
0.038), while arm~N realized 0.982 against 0.975 promised (Chicago, 2{,}570 answered) and 0.929 against 0.910 (OSHA, 294
answered). Arm~N's difference is at or below zero in every Chicago range and in every OSHA range where it answered; in
the first OSHA range it answered nothing, so there the comparison is between a keeper that abstained and one that
answered and was wrong. Further on, arm~C recovers in Chicago (0.972 realized against 0.967 promised over records
20{,}000--120{,}000) and not in OSHA within the 17{,}746 records available.

\paragraph{Detection.} Table~\ref{tab:watch}.

\begin{table}[h]
\centering
\caption{The realized-precision watch, not told about the revision.}
\label{tab:watch}
\begin{tabular}{lrrrrr}
\toprule
Stream & Records followed & Revision at & Alarms & Before the revision & First alarm after \\
\midrule
Chicago & 210{,}842 & 90{,}842 & 2 & 0 & 54 records \\
OSHA & 105{,}991 & 88{,}245 & 3 & 0 & 163 records \\
\bottomrule
\end{tabular}
\end{table}

The later alarms are the carried-over memory still answering with old codes more than 5{,}000 records on; the run does
nothing with an alarm. On thirteen public streams without a revision the same watch with the same constants raised no
alarm. On three further streams not used in setting it, it raised one: a nine-day run of recalls from one firm, all
answered wrongly by the memory --- a burst from one source, not a revision. For this reason the watch only reports; we
do not restart the memory on an alarm.

\paragraph{Frozen thresholds.} On the pre-revision hold-out every frozen policy met its target (0.949--0.952 precision
on Chicago, 0.901--0.913 on OSHA). Table~\ref{tab:frozen} gives the first 1{,}000 records after the revision.

\begin{table}[h]
\centering
\caption{Score-threshold policies after the revision.}
\label{tab:frozen}
\begin{tabular}{llrr}
\toprule
Stream & Policy & Share answered (first 1{,}000) & Precision \\
\midrule
Chicago & CalibRoute, frozen & 0.994 & 0.004 \\
Chicago & MAPIE, frozen & 1.000 & 0.004 \\
Chicago & CalibRoute with its shift check & 0 (all 480 batches flagged) & -- \\
Chicago & retrained and re-thresholded & 0.353 & 0.941 \\
OSHA & CalibRoute, frozen & 0.947 & 0.064 \\
OSHA & MAPIE, frozen & 0.938 & 0.064 \\
OSHA & CalibRoute with its shift check & 0.947 (no batch flagged) & 0.064 \\
OSHA & retrained and re-thresholded & 0 & -- \\
\bottomrule
\end{tabular}
\end{table}

The shift check stopped the frozen policy where the classifier's confidences moved (Chicago) and not where they did not
(OSHA): the classifier stayed as confident as before while the codes it pointed to had been renumbered. The two ways of
choosing again after the revision --- the keeper started fresh and the retrained, re-thresholded classifier --- both
stayed near the target, with different coverage (records 1{,}000--5{,}000: keeper 0.64 at 0.98 and 0.07 at 0.93;
retrained 0.86 at 0.89 and 0.09 at 0.92, Chicago and OSHA respectively).

\section{What this does not show}

\begin{itemize}
  \item \textbf{Two streams, one kind of revision.} Both revisions rewrite code numbers. Nothing here is about a
  revision that changes no code, or about other streams.
  \item \textbf{Derived settlements.} A record is settled with the code the source lists. In use, people may accept
  wrong answers or correct right ones, and only some answers may be checked; the watch is as good as the settlements it
  gets.
  \item \textbf{One setting.} Target 0.9, a window of 100 and a margin of 0.15. The watch assumes the default threshold
  rule, which keeps a floor under every band of evidence; with a rule that pools all bands it raised 21 false alarms
  across the thirteen streams without a revision, against none.
  \item \textbf{Splitting has a price.} Starting the memory at a boundary where codes did not change costs answers on
  attributes with many values. The conclusion is not ``always start afresh'' but ``an edition is a unit of memory: split
  it where an edition is declared''.
  \item \textbf{Not an argument against retraining.} A retrained threshold kept close to its target too; these runs
  show what freezing a policy across a revision costs, not that a keeper does what retraining cannot.
  \item \textbf{Small first ranges.} The first range is 1{,}000 records; on OSHA arm~C answered 153 of them.
\end{itemize}

\section{The request}

We would like to hear of any public stream where (a) a classification edition rewrites code numbers and a memory
carried across it still keeps a stated precision in the first records after it, or (b) the realized-precision watch
stays silent for more than 1{,}000 records after such a revision, or (c) it fires on a stream without a revision for a
reason other than a burst from one source. Each would narrow or break a claim above. The kit~\cite{kit} builds both
streams from the public data and prints every table in this note with \texttt{iyulab-gil} 0.10.0 from PyPI:

\begin{verbatim}
uv run --no-project python build_world.py chicago inspections-2017-.csv \
  --out chicago.stream.jsonl.gz
uv run --no-project --isolated --with iyulab-gil==0.10.0 python carryover.py \
  chicago.stream.jsonl.gz --slot violation --boundary 2018 --after 120000 \
  --arms N,C --out chicago.result.json
\end{verbatim}

The Chicago portal changes continuously; a later download reproduces the pattern and every rate within 0.002, not every
last digit. The OSHA run is fixed by the archive's SHA-256 and prints its tables exactly.

\begin{thebibliography}{9}
\bibitem{geifman2017} Y.~Geifman, R.~El-Yaniv. Selective classification for deep neural networks. NeurIPS 2017.
arXiv:1705.08500.
\bibitem{ltt} A.~N. Angelopoulos, S.~Bates, E.~J. Cand\`es, M.~I. Jordan, L.~Lei. Learn then Test: calibrating
predictive algorithms to achieve risk control. Annals of Applied Statistics; arXiv:2110.01052.
\bibitem{crc} A.~N. Angelopoulos, S.~Bates, A.~Fisch, L.~Lei, T.~Schuster. Conformal risk control. ICLR 2024;
arXiv:2208.02814.
\bibitem{gil} Gil. \url{https://github.com/iyulab/gil} (package \texttt{iyulab-gil} on PyPI and crates.io, \texttt{Gil}
on NuGet).
\bibitem{calibroute} CalibRoute, \texttt{calibroute-ai} 0.4.0. \url{https://pypi.org/project/calibroute-ai/}
\bibitem{mapie} MAPIE 1.6.0. \url{https://pypi.org/project/mapie/}
\bibitem{kit} Gil replication kits. \url{https://github.com/iyulab/gil-replication}
\end{thebibliography}

\end{document}
